\section{Lingkungan Implementasi}
\label{sec:lingkungan-implementasi}

Implementasi sistem dilakukan pada lingkungan prototipe yang merepresentasikan
komponen utama DecMed, yaitu \textit{Patient Client}, \textit{Hospital Client},
\textit{Ministry Client}, PRE Server, penyimpanan \textit{off-chain} IPFS
Cluster, IOTA Gas Station, \textit{smart contract} IOTA Move, dan layanan
Audit Log Service (ALS) yang berjalan pada jaringan IOTA Localnet. Lingkungan
ini dipilih agar rancangan pada Bab III dapat diuji tanpa bergantung pada
integrasi langsung dengan sistem informasi rumah sakit yang sesungguhnya.

Komponen ALS dan PRE Server dijalankan bersama pada lingkungan lokal bersama
ketiga aplikasi \textit{client}. Komponen pendukung yang membutuhkan ketersediaan
jaringan, yaitu IPFS Cluster dan IOTA Gas Station, dijalankan pada Virtual
Private Server (VPS). \textit{Smart contract} IOTA Move di-\textit{deploy} pada
jaringan IOTA Localnet yang juga berjalan di VPS. IOTA Localnet digunakan
karena sudah memadai untuk menyimulasikan interaksi dengan jaringan IOTA.

Spesifikasi perangkat lokal yang digunakan adalah sebagai berikut.
\begin{enumerate}
    \item Sistem Operasi: Ubuntu 24.04.4, berjalan di lingkungan WSL2.
    \item CPU: Intel Core i7-12700H (6 \textit{core} fisik dan 12 \textit{logical thread} dialokasikan untuk WSL2).
    \item Memori RAM: 32 GB (16 GB teralokasi untuk WSL2).
    \item Perangkat bantu pengembangan: Rust (\texttt{rustc 1.85.0},
    \texttt{cargo 1.85.0}), Node.js v22.14.0, dan Docker 27.4.0.
\end{enumerate}

Sementara itu, spesifikasi VPS yang digunakan adalah sebagai berikut.
\begin{enumerate}
    \item Sistem Operasi: Ubuntu 24.04.4 LTS.
    \item Jumlah vCPU: 4 vCPU.
    \item Memori RAM: 8 GB.
    \item Perangkat bantu: Docker 27.4.0.
\end{enumerate}

\section{Implementasi Sistem}

Bagian ini menjabarkan realisasi teknis dari rancangan pada Bab III. Narasi
konsep dan alur yang sudah ditetapkan di Bab III dirujuk ulang; fokus di sini
adalah komponen kode, antarmuka, dan perilaku konkret sistem.

Daftar komponen dan keterkaitannya dengan rancangan ditunjukkan pada
Tabel~\ref{tab:daftar-implementasi-als}. ID I-ALS dipakai sebagai rujukan.
R-TA-01 (rancangan arsitektur solusi) mencakup keseluruhan sistem sehingga
tidak muncul sebagai baris tersendiri.

\begin{table}[htbp]
\centering
\caption{Daftar Implementasi Sistem Audit Log}
\label{tab:daftar-implementasi-als}
\renewcommand{\arraystretch}{1.3}
\begin{tabular}{|p{1.8cm}|p{5cm}|p{5.5cm}|}
\hline
\textbf{ID} & \textbf{Komponen} & \textbf{Merealisasikan Rancangan} \\
\hline
I-ALS-01 & Skema tipe data \texttt{AuditEvent}, \texttt{Event}, \texttt{AuditEventDetails}, \texttt{AuditRecord}, dan \texttt{AuditBatch} beserta enum pendukungnya. & Katalog \textit{audit event} (EV1--EV13) dan skema atribut audit pada Subbab~\ref{subsec:rancangan-audit-record}. \\
\hline
I-ALS-02 & Tipe data \texttt{EncryptedSignedEvent} dan mekanisme enkripsi-tanda tangan berlapis pada sisi \textit{event source}. & Subbab~\ref{subsec:rancangan-pembangkitan-event} (lapisan enkapsulasi \textit{event}). \\
\hline
I-ALS-03 & \textit{Endpoint} \texttt{POST /api/events} dan mekanisme dekripsi serta verifikasi \textit{encrypted signed event} pada ALS. & Subbab~\ref{subsec:rancangan-pembangkitan-event} (komponen Audit Receiver). \\
\hline
I-ALS-04 & Komponen \texttt{AuditLogger} dengan antrian asinkron internal. & Subbab~\ref{subsec:rancangan-pembentukan-record} (komponen Audit Logger dan mekanisme \textit{hash-chaining}). \\
\hline
I-ALS-05 & \textit{Worker} rotasi log, unggah IPFS, dan publikasi metadata ke IOTA. & Subbab~\ref{subsec:rancangan-rotasi-log}--\ref{subsec:rancangan-publikasi-iota} (Log Rotation \& Archival Worker). \\
\hline
I-ALS-06 & Modul \texttt{als} dengan \texttt{ALSClient}, \texttt{AlsQueue}, dan \textit{retry worker} pada setiap \textit{event source}. & Subbab~\ref{subsec:rancangan-pembangkitan-event} dan \ref{subsubsec:rancangan-mekanisme-pengiriman} (integrasi ALS ke dalam arsitektur DecMed). \\
\hline
\end{tabular}
\end{table}

Tabel~\ref{tab:daftar-implementasi-als} menjadi acuan urutan pembahasan
implementasi berikut, dari skema tipe data, mekanisme enkapsulasi \textit{event},
penerimaan dan verifikasi oleh ALS, pembentukan \textit{audit record} dengan
\textit{hash-chaining}, rotasi log dan pengarsipan, hingga integrasi ke seluruh
\textit{event source} DecMed.

\subsection{Batasan Implementasi}

Implementasi mempertahankan arsitektur kerangka kerja DecMed sebagaimana
dijelaskan pada Subbab~\ref{subsec:rancangan-arsitektur}. Komponen perangkat
lunak ALS diimplementasikan sebagai \textit{crate} Rust yang berdiri sendiri.
Dependensi utama yang digunakan ditunjukkan pada Tabel~\ref{tab:dependensi-als}.
Batasan cakupan implementasi, konfigurasi IPFS \textit{pinning}, dan nilai
\texttt{LOG\_ROTATION\_INTERVAL\_SECS} diuraikan pada
Subbab~\ref{sec:pengujian}.

\begin{table}[htbp]
\centering
\caption{Dependensi Utama Audit Log Service}
\label{tab:dependensi-als}
\renewcommand{\arraystretch}{1.3}
\begin{tabular}{|p{3.5cm}|p{2cm}|p{6.5cm}|}
\hline
\textbf{Dependensi} & \textbf{Versi} & \textbf{Fungsi} \\
\hline
\texttt{axum} & 0.8.4 & \textit{Web framework} HTTP untuk \textit{endpoint} penerimaan \textit{audit event}. \\
\hline
\texttt{tokio} & 1.44.2 & \textit{Async runtime} Rust untuk pemrosesan konkuren. \\
\hline
\texttt{serde} / \texttt{serde\_json} & 1.x & Serialisasi dan deserialisasi data JSON. \\
\hline
\texttt{aes-gcm} & 0.10 & Enkripsi dan dekripsi AES-256-GCM pada lapisan transport \textit{event}. \\
\hline
\texttt{p256} & 0.13 & ECDH P-256 untuk skema ECIES pembungkusan kunci AES. \\
\hline
\texttt{iota-sdk} / \texttt{iota-types} & 1.x & Tipe kriptografi IOTA (\texttt{IotaKeyPair}, \texttt{IotaSignature}) dan interaksi jaringan IOTA. \\
\hline
\texttt{shared-crypto} & 1.x & Pembungkusan \textit{intent message} IOTA untuk penandatanganan \textit{event}. \\
\hline
\texttt{sha2} & 0.10 & Komputasi \textit{hash} SHA-256 untuk mekanisme \textit{hash-chaining} dan KDF ECIES. \\
\hline
\texttt{uuid} & 1.x & Pembangkitan \texttt{record\_id} berbasis UUID v7 yang bersifat monoton. \\
\hline
\texttt{chrono} & 0.4.39 & Penanganan waktu dan \textit{timestamp} dalam format UTC. \\
\hline
\texttt{reqwest} & 0.12 & Klien HTTP untuk unggah berkas ke IPFS dan pengiriman \textit{event} ke ALS. \\
\hline
\texttt{tokio-util} & 0.7.11 & \textit{Streaming} berkas saat unggah ke IPFS menggunakan \texttt{BytesCodec}. \\
\hline
\end{tabular}
\end{table}

\subsection{Implementasi Skema Tipe Data Audit}
\label{subsec:impl-skema-tipe-data}

Implementasi skema tipe data audit merealisasikan katalog \textit{audit event}
dan skema atribut yang telah ditentukan pada Subbab~\ref{subsec:rancangan-audit-record}.
Tipe-tipe data pada sisi ALS didefinisikan pada berkas
\texttt{audit-trail/src/types.rs}, sedangkan tipe-tipe setara yang digunakan
pada setiap \textit{event source} didefinisikan masing-masing pada berkas
\texttt{als/types.rs} milik komponen bersangkutan. Tipe utama yang menyusun
skema ini adalah \texttt{EncryptedSignedEvent}, \texttt{AuditEvent}, \texttt{Event},
empat \textit{enum} pendukung (\texttt{AuditOutcome}, \texttt{AuditSourceComponent},
\texttt{AuditActionType}, \texttt{AuditActorType}, \texttt{AuditTargetObjectType}),
\texttt{AuditEventDetails}, \texttt{AuditRecord}, dan \texttt{AuditBatch}.

\subsubsection{Tipe Data EncryptedSignedEvent}

\texttt{EncryptedSignedEvent} adalah struktur yang dikirimkan oleh setiap
\textit{event source} ke ALS melalui \textit{endpoint} HTTP. Tipe ini mewakili
\textit{event} yang telah dienkripsi dan ditandatangani, sehingga isi
\textit{event} tidak dapat dibaca oleh pihak ketiga yang menyadap saluran
komunikasi. Tipe ini memuat lima bidang: \texttt{enc\_aes\_key} sebagai kunci
sesi AES yang telah dibungkus menggunakan ECIES P-256 dalam representasi Base64,
\texttt{ciphertext} sebagai \textit{payload audit event} yang telah dienkripsi
AES-256-GCM dalam representasi Base64, \texttt{nonce} sebagai \textit{nonce}
AES-GCM sepanjang 12 \textit{byte} dalam representasi Base64, \texttt{signature}
sebagai tanda tangan \texttt{IotaSignature} atas \texttt{ciphertext} dalam
representasi Base64, dan \texttt{iota\_address} sebagai alamat IOTA pengirim
yang digunakan ALS untuk memverifikasi tanda tangan. Implementasi
\texttt{EncryptedSignedEvent} ditunjukkan pada Kode~\ref{code:encrypted-signed-event}.

\begin{listing}[ht]
\caption{Implementasi \texttt{EncryptedSignedEvent}}
\label{code:encrypted-signed-event}
\begin{minted}[linenos, breaklines]{rust}
#[derive(Debug, Clone, Serialize, Deserialize)]
pub struct EncryptedSignedEvent {
    pub enc_aes_key: String,  // base64: ECIES(ephemeral_pubkey.wrap_nonce.enc_key)
    pub ciphertext: String,   // base64: AuditEvent terenkripsi AES-256-GCM
    pub nonce: String,        // base64: nonce AES-GCM (12 bytes)
    pub signature: String,    // base64: IotaSignature atas ciphertext bytes
    pub iota_address: String, // untuk verifikasi signature di ALS
}
\end{minted}
\end{listing}

Perlu dicatat bahwa pada sisi ALS, terdapat juga tipe \texttt{SignedEvent}
yang digunakan secara internal untuk keperluan pengujian langsung ke
\textit{endpoint}, memuat bidang \texttt{payload} (JSON \textit{plaintext}),
\texttt{signature} (heksadesimal Ed25519), dan \texttt{public\_key}
(heksadesimal). Pada alur produksi, \textit{event source} selalu mengirimkan
\texttt{EncryptedSignedEvent}, bukan \texttt{SignedEvent}.

\subsubsection{Tipe Data AuditEvent dan Event}

\texttt{AuditEvent} mewakili satu peristiwa audit yang telah diekstrak dan
terdekripsi dari \texttt{EncryptedSignedEvent}. Tipe ini memuat tiga bidang:
\texttt{source\_component} bertipe \texttt{AuditSourceComponent} (komponen
sumber pembangkit \textit{event}), \texttt{source\_timestamp} bertipe
\texttt{DateTime<Utc>} (cap waktu dari sisi \textit{event source} pada saat
kejadian berlangsung), dan \texttt{event} bertipe \texttt{Event} yang memuat
seluruh atribut detail \textit{event}.

Pemisahan \texttt{source\_timestamp} dari \texttt{timestamp} pada
\texttt{AuditRecord} dilakukan secara sengaja: \texttt{source\_timestamp}
mencatat waktu kejadian menurut jam \textit{event source}, sedangkan
\texttt{timestamp} pada \texttt{AuditRecord} mencatat waktu pembentukan
\textit{record} menurut jam ALS. Kedua nilai ini dapat berbeda apabila terdapat
latensi jaringan atau antrean, sehingga keduanya diperlukan untuk keperluan
rekonstruksi urutan kejadian.

\texttt{Event} memuat tujuh bidang yang merealisasikan skema atribut umum dan
atribut detail: \texttt{actor\_id} (identitas aktor yang memicu aktivitas),
\texttt{actor\_type} bertipe \texttt{AuditActorType} (kategori aktor),
\texttt{target\_object\_type} bertipe \texttt{AuditTargetObjectType} (kategori
objek target), \texttt{target\_object} (identitas spesifik objek target),
\texttt{outcome} bertipe \texttt{AuditOutcome} (hasil aktivitas),
\texttt{action\_type} bertipe \texttt{AuditActionType} (jenis tindakan), dan
\texttt{details} bertipe \texttt{AuditEventDetails} yang diserialisasi sejajar
menggunakan anotasi \texttt{\#[serde(flatten)]}. Implementasi kedua struct
ditunjukkan pada Kode~\ref{code:audit-event}.

\begin{listing}[ht]
\caption{Implementasi \texttt{AuditEvent} dan \texttt{Event}}
\label{code:audit-event}
\begin{minted}[linenos, breaklines]{rust}
#[derive(Clone, Debug, Deserialize, Serialize)]
pub struct AuditEvent {
    pub source_component: AuditSourceComponent,
    pub source_timestamp: DateTime<Utc>,
    pub event: Event,
}

#[derive(Clone, Debug, Deserialize, Serialize)]
pub struct Event {
    pub actor_id: String,
    pub actor_type: AuditActorType,
    pub target_object_type: AuditTargetObjectType,
    pub target_object: String,
    pub outcome: AuditOutcome,
    pub action_type: AuditActionType,

    #[serde(flatten)]
    pub details: AuditEventDetails,
}
\end{minted}
\end{listing}

\subsubsection{Enum AuditOutcome dan AuditSourceComponent}

\texttt{AuditOutcome} diimplementasikan sebagai \textit{enum} dengan empat
varian: \texttt{Success}, \texttt{Failure}, \texttt{Denied}, dan
\texttt{Unknown}, merealisasikan atribut \texttt{outcome} pada skema atribut
umum. Anotasi \texttt{\#[serde(rename\_all = "PascalCase")]} memastikan nilai
JSON berupa \texttt{"Success"}, \texttt{"Failure"}, dan seterusnya.

\texttt{AuditSourceComponent} diimplementasikan sebagai \textit{enum} dengan
empat varian yang merealisasikan seluruh \textit{event source} DecMed:
\texttt{HospitalClient}, \texttt{PatientClient}, \texttt{MinistryClient}, dan
\texttt{ProxyReencryption}. Anotasi \texttt{\#[serde(rename\_all = "PascalCase")]}
juga diterapkan pada \textit{enum} ini. Penggunaan \textit{enum} berketik
(\textit{typed enum}) menggantikan representasi \texttt{String} memastikan
bahwa hanya nilai yang valid secara semantik yang dapat dimasukkan ke bidang
\texttt{source\_component}; nilai yang tidak dikenal akan ditolak pada tahap
deserialisasi. Implementasi kedua \textit{enum} ini ditunjukkan pada
Kode~\ref{code:audit-outcome-source}.

\begin{listing}[ht]
\caption{Implementasi \texttt{AuditOutcome} dan \texttt{AuditSourceComponent}}
\label{code:audit-outcome-source}
\begin{minted}[linenos, breaklines]{rust}
#[derive(Clone, Debug, Deserialize, Serialize)]
#[serde(rename_all = "PascalCase")]
pub enum AuditOutcome {
    Success,
    Failure,
    Denied,
    Unknown,
}

#[derive(Clone, Debug, Deserialize, Serialize)]
#[serde(rename_all = "PascalCase")]
pub enum AuditSourceComponent {
    HospitalClient,
    PatientClient,
    MinistryClient,
    ProxyReencryption,
}
\end{minted}
\end{listing}

\subsubsection{Enum AuditActionType}

\texttt{AuditActionType} diimplementasikan sebagai \textit{enum} yang
mengkategorikan jenis tindakan yang dilakukan oleh aktor. Varian yang
didefinisikan adalah \texttt{Create}, \texttt{Read}, \texttt{Update},
\texttt{Delete}, \texttt{Validate}, dan \texttt{Execute}. Pada \textit{Hospital
Client} dan \textit{Patient Client}, varian tambahan \texttt{SignIn} ditambahkan
untuk mencatat peristiwa autentikasi sesi.

Anotasi \texttt{\#[serde(rename\_all = "SCREAMING\_SNAKE\_CASE")]} diterapkan
sehingga nilai JSON berupa \texttt{"CREATE"}, \texttt{"READ"}, \texttt{"SIGN\_IN"},
dan seterusnya. Konsistensi format ini mempermudah pemrosesan log secara
otomatis karena setiap nilai \texttt{action\_type} dapat dicocokkan secara tepat
terhadap daftar yang telah ditentukan. Implementasinya ditunjukkan pada
Kode~\ref{code:audit-action-type}.

\begin{listing}[ht]
\caption{Implementasi \texttt{AuditActionType}}
\label{code:audit-action-type}
\begin{minted}[linenos, breaklines]{rust}
// Versi Hospital Client dan Patient Client (dengan SignIn)
#[derive(Clone, Debug, Deserialize, Serialize)]
#[serde(rename_all = "SCREAMING_SNAKE_CASE")]
pub enum AuditActionType {
    Create,
    Read,
    Update,
    Delete,
    Validate,
    Execute,
    SignIn,
}
\end{minted}
\end{listing}

\subsubsection{Enum AuditActorType dan AuditTargetObjectType}

\texttt{AuditActorType} mengkategorikan jenis aktor yang memicu \textit{event}.
Varian yang didefinisikan mencakup \texttt{AdministrativePersonnel},
\texttt{MedicalPersonnel}, \texttt{Patient}, \texttt{Admin}, \texttt{Ministry},
dan \texttt{PREServer}. Pada \textit{Hospital Client} dan \textit{Patient Client},
dua varian tambahan \texttt{Unknown} dan \texttt{Client} disertakan untuk
menangani skenario di mana identitas aktor belum dapat ditentukan saat
\textit{event} dibangkitkan.

Untuk memudahkan integrasi dengan sistem otorisasi yang sudah ada, setiap
komponen mengimplementasikan \textit{trait} konversi \texttt{From} dari tipe
peran lokal ke \texttt{AuditActorType}. Pada PRE Server,
\texttt{From<AuthRole>} mengonversi \texttt{AuthRole::AdministrativePersonnel},
\texttt{AuthRole::MedicalPersonnel}, dan \texttt{AuthRole::Patient} ke varian
\texttt{AuditActorType} yang bersesuaian. Pada \textit{Hospital Client},
\texttt{From<HospitalPersonnelRole>} mengonversi \texttt{HospitalPersonnelRole::Admin},
\texttt{HospitalPersonnelRole::AdministrativePersonnel}, dan
\texttt{HospitalPersonnelRole::MedicalPersonnel}. Selain itu, \textit{trait}
\texttt{FromStr} juga diimplementasikan pada \textit{Hospital Client} dan
\textit{Patient Client} untuk mengonversi nilai \textit{string} (misalnya
dari sesi yang tersimpan) ke \texttt{AuditActorType}. Implementasi konversi
pada PRE Server ditunjukkan pada Kode~\ref{code:actor-type-from}.

\begin{listing}[ht]
\caption{Implementasi \texttt{AuditActorType} dan konversi \texttt{From<AuthRole>} pada PRE Server}
\label{code:actor-type-from}
\begin{minted}[linenos, breaklines]{rust}
#[derive(Clone, Debug, Deserialize, Serialize)]
pub enum AuditActorType {
    AdministrativePersonnel,
    MedicalPersonnel,
    Patient,
    Admin,
    Ministry,
    PREServer,
}

impl From<AuthRole> for AuditActorType {
    fn from(role: AuthRole) -> Self {
        match role {
            AuthRole::AdministrativePersonnel =>
                AuditActorType::AdministrativePersonnel,
            AuthRole::MedicalPersonnel =>
                AuditActorType::MedicalPersonnel,
            AuthRole::Patient =>
                AuditActorType::Patient,
        }
    }
}
\end{minted}
\end{listing}

\texttt{AuditTargetObjectType} mengkategorikan jenis objek yang menjadi target
aktivitas. Varian yang digunakan bervariasi antar \textit{event source} sesuai
objek yang dapat diakses oleh masing-masing komponen. PRE Server mendefinisikan
varian yang lebih lengkap, mencakup \texttt{ActivationKey}, \texttt{AccessCapability},
\texttt{MedicalRecord}, \texttt{MedicalRecordMetadata}, \texttt{AccessDelegationQR},
\texttt{AdministrativeData}, \texttt{Nonce}, \texttt{AccessKeys},
\texttt{Transaction}, \texttt{IPFSObject}, \texttt{GasReservation},
\texttt{PREEndPoint}, dan \texttt{OnChainData}. \textit{Client} aplikasi
mendefinisikan subset dari varian tersebut yang relevan dengan operasi yang
dapat dipicu dari antarmuka pengguna.

\subsubsection{Tipe Data AuditEventDetails}

\texttt{AuditEventDetails} diimplementasikan sebagai \textit{tagged enum} dengan
anotasi \texttt{\#[serde(tag = "event\_type")]} sehingga nilai \texttt{event\_type}
pada JSON yang masuk menentukan varian \textit{enum} yang akan di-\textit{parse}.
Terdapat 13 varian yang masing-masing merealisasikan satu jenis \textit{audit
event} (EV1--EV13) sesuai Tabel~\ref{tab:threat-event-mapping}.

Perbedaan antarjenis \textit{event} terletak seluruhnya pada bidang dalam varian
masing-masing. Sebagai contoh, varian \texttt{EV7} (\textit{PRE Service
Operation}) memuat \texttt{endpoint\_called}, \texttt{request\_id},
\texttt{caller\_component}, \texttt{channel\_encryption}, dan
\texttt{jwt\_purpose}. Pada \textit{client} aplikasi, varian \texttt{EV7}
ditambahkan dua bidang observabilitas yaitu \texttt{http\_status\_code}
(bertipe \texttt{u16}) dan \texttt{latency\_ms} (bertipe \texttt{u64}) yang
tidak didefinisikan pada sisi PRE Server karena PRE Server mengukur latensi
dari perspektif yang berbeda. Varian \texttt{EV13} (validasi \textit{capability})
memuat \texttt{capability\_id}, \texttt{requester\_id},
\texttt{validation\_result} (bertipe \texttt{bool}), dan
\texttt{rejection\_reason} (bertipe \texttt{Option<String>} yang bernilai
\texttt{None} apabila validasi berhasil). Implementasi seluruh varian
\texttt{AuditEventDetails} ditunjukkan pada Lampiran~\ref{code:audit-event-details}.

\subsubsection{Tipe Data AuditRecord}

\texttt{AuditRecord} mewakili catatan audit yang telah terbentuk dan siap
disimpan ke berkas log. Tipe ini memuat empat bidang tambahan di luar konten
\textit{event}: \texttt{record\_id} bertipe \texttt{Uuid} (UUID v7 yang
mengandung informasi waktu sehingga bersifat monoton dan dapat diurutkan secara
kronologis), \texttt{timestamp} bertipe \texttt{DateTime<Utc>} (waktu pembentukan
\textit{record} menurut jam ALS), \texttt{prev\_record\_hash} bertipe
\texttt{Option<String>} (nilai \textit{hash} SHA-256 dari \textit{record}
sebelumnya, atau \texttt{None} apabila merupakan \textit{record} pertama dalam
rantai), dan \texttt{record\_hash} bertipe \texttt{String} (nilai \textit{hash}
SHA-256 yang dihitung dari kombinasi keempat bidang tersebut bersama isi
\textit{event}). Seluruh atribut \texttt{AuditEvent} diserialisasi sejajar ke
dalam \texttt{AuditRecord} menggunakan \texttt{\#[serde(flatten)]}, sehingga
berkas JSON-Lines yang dihasilkan memuat seluruh atribut pada satu tingkat yang
sama. Implementasi \texttt{AuditRecord} ditunjukkan pada Kode~\ref{code:audit-record}.

\begin{minted}{rust}
#[derive(Debug, Clone, Serialize, Deserialize)]
pub struct AuditRecord {
    pub record_id: Uuid,
    pub timestamp: DateTime<Utc>,
    pub prev_record_hash: Option<String>,
    pub record_hash: String,

    #[serde(flatten)]
    pub event: AuditEvent,
}
\end{minted}

\captionof{listing}{Implementasi \texttt{AuditRecord}}
\label{code:audit-record}

\subsubsection{Tipe Data AuditBatch}

\texttt{AuditBatch} merepresentasikan ringkasan satu berkas log yang telah
dirotasi dan diarsipkan ke IPFS. Tipe ini tidak disimpan dalam berkas log itu
sendiri, melainkan dibentuk oleh \textit{worker} setelah proses rotasi selesai
sebagai metadata yang merangkum isi berkas arsip. Bidang-bidang pada
\texttt{AuditBatch} mencakup identitas dan rentang waktu \textit{batch}
(\texttt{batch\_id}, \texttt{start\_time}, \texttt{end\_time}), ringkasan isi
(\texttt{record\_count}, \texttt{first\_record\_id}, \texttt{last\_record\_id}),
nilai \textit{hash} rekam pertama dan terakhir untuk verifikasi ujung rantai
(\texttt{first\_record\_hash}, \texttt{final\_record\_hash}), serta informasi
persistensi (\texttt{file\_hash} sebagai \textit{hash} SHA-256 berkas arsip
secara keseluruhan, dan \texttt{ipfs\_cid} sebagai CID yang dihasilkan setelah
unggah berhasil). Implementasi \texttt{AuditBatch} ditunjukkan pada
Kode~\ref{code:audit-batch}.

\begin{minted}{rust}
#[derive(Debug, Clone, Serialize, Deserialize)]
pub struct AuditBatch {
    pub batch_id: Uuid,
    pub start_time: DateTime<Utc>,
    pub end_time: DateTime<Utc>,
    pub record_count: u64,
    pub first_record_id: Uuid,
    pub last_record_id: Uuid,
    pub first_record_hash: String,
    pub final_record_hash: String,
    pub file_hash: String,
    pub ipfs_cid: String,
}
\end{minted}

\captionof{listing}{Implementasi \texttt{AuditBatch}}
\label{code:audit-batch}

\subsection{Implementasi Mekanisme Enkapsulasi Event pada Event Source}
\label{subsec:impl-enkapsulasi-event}

Sebelum \textit{event} dikirimkan ke ALS, setiap \textit{event source}
mengenkapsulasi \textit{event} melalui tujuh langkah berurutan yang
diimplementasikan pada metode privat \texttt{build\_and\_send} milik
\texttt{ALSClient}. Langkah-langkah tersebut ditunjukkan pada
Kode~\ref{code:build-and-send-steps}.

\begin{listing}[ht]
\caption{Alur enkapsulasi \textit{event} pada \texttt{ALSClient::build\_and\_send}}
\label{code:build-and-send-steps}
\begin{minted}[linenos, breaklines]{rust}
// Step 1: Serialize AuditEvent ke bytes JSON
let payload_bytes = serde_json::to_vec(&event)?;

// Step 2: Enkripsi payload dengan AES-256-GCM (kunci sesi acak)
let encrypted = aes_encrypt(&payload_bytes)?;

// Step 3: Bungkus kunci AES dengan public key ALS menggunakan ECIES P-256
let enc_aes_key = ecies_encrypt_key(&encrypted.key, ALS_SERVER_PUBLIC_KEY)?;

// Step 4: Tandatangani ciphertext dengan IotaKeyPair (Encrypt-then-Sign)
let intent_msg = IntentMessage::new(
    Intent::personal_message(),
    encrypted.ciphertext.clone(),
);
let signature = IotaSignature::new_secure(&intent_msg, iota_key_pair);

// Step 5: Bangun EncryptedSignedEvent
let signed_event = EncryptedSignedEvent {
    enc_aes_key,
    ciphertext: STANDARD.encode(&encrypted.ciphertext),
    nonce: STANDARD.encode(&encrypted.nonce),
    signature: signature.encode_base64(),
    iota_address: iota_address.clone(),
};

// Step 6: Simpan ke AlsQueue terlebih dahulu (durabilitas)
AlsQueue::push(new_queue_entry(signed_json.clone(), label)).await?;

// Step 7: Coba kirim langsung; jika gagal, retry worker akan menangani
reqwest::Client::new()
    .post(ALS_ENDPOINT)
    .header("Content-Type", "application/json")
    .body(signed_json)
    .send()
    .await;
\end{minted}
\end{listing}

Langkah 1 hingga 5 merealisasikan mekanisme enkapsulasi berlapis yang dirancang
pada Subbab~\ref{subsubsec:rancangan-enkapsulasi-event}. Langkah 2 dan 3
bersama-sama merealisasikan lapisan enkripsi: \textit{payload} dienkripsi
dengan AES-256-GCM menggunakan kunci sesi acak, dan kunci sesi tersebut
kemudian dibungkus dengan kunci publik P-256 ALS melalui skema ECIES.
Implementasi ECIES pada fungsi \texttt{ecies\_encrypt\_key} (berkas
\texttt{als/crypto.rs}) menggunakan \textit{ephemeral keypair} P-256 yang
dibangkitkan baru pada setiap pemanggilan, melakukan ECDH terhadap kunci publik
ALS, menerapkan SHA-256 sebagai KDF (\textit{key derivation function}) atas
\textit{shared secret}, kemudian mengenkripsi kunci AES dengan kunci turunan
tersebut menggunakan AES-256-GCM. Hasilnya dikodekan sebagai satu \textit{string}
Base64 yang menggabungkan \textit{ephemeral public key}, \textit{wrap nonce},
dan \textit{encrypted AES key}.

Langkah 4 merealisasikan lapisan tanda tangan menggunakan pola
\textit{Encrypt-then-Sign}: yang ditandatangani adalah \textit{ciphertext},
bukan \textit{plaintext}. Penandatanganan menggunakan \texttt{IotaKeyPair}
milik aktor yang terlibat pada saat kejadian berlangsung, yang diperoleh dari
\texttt{AppState} komponen. Penggunaan kunci IOTA alih-alih kunci Ed25519
terpisah memiliki implikasi penting: kunci tersebut sudah terikat pada identitas
aktor yang terdaftar pada \textit{ledger} IOTA, sehingga \textit{non-repudiation}
dapat ditegakkan secara kriptografis terhadap identitas yang terregistrasi pada
sistem.

Langkah 6 merupakan mekanisme durabilitas yang dibahas lebih lanjut pada
Subbab~\ref{subsec:impl-als-queue}.

\subsection{Implementasi AlsQueue dan Retry Worker}
\label{subsec:impl-als-queue}

Untuk menjamin tidak ada \textit{event} yang hilang akibat kegagalan jaringan
sementara, setiap \textit{event source} mengimplementasikan \texttt{AlsQueue},
sebuah antrian berbasis berkas (\textit{file-based queue}) yang menyimpan
\texttt{EncryptedSignedEvent} yang belum berhasil dikirimkan. Implementasi
\texttt{AlsQueue} terdapat pada berkas \texttt{als/queue.rs} di setiap
\textit{event source}.

Setiap entri antrian direpresentasikan sebagai \texttt{QueueEntry} yang memuat
\texttt{id} (UUID v4), \texttt{signed\_payload} (JSON \texttt{EncryptedSignedEvent}),
\texttt{label} (nama \textit{event} untuk pencatatan diagnostik),
\texttt{created\_at} (waktu pembuatan entri), dan \texttt{attempt\_count}
(jumlah percobaan pengiriman yang telah dilakukan). Antrian disimpan sebagai
berkas JSON-Lines pada direktori lokal sehingga entri yang belum terkirim tetap
bertahan meskipun aplikasi diulang.

Pola pengiriman yang diterapkan adalah \textit{store-then-send}: entri disimpan
ke antrian terlebih dahulu (Langkah 6 pada Kode~\ref{code:build-and-send-steps}),
kemudian pengiriman langsung dicoba (Langkah 7). Apabila pengiriman langsung
berhasil, entri dihapus dari antrian. Apabila gagal, entri tetap berada di
antrian dan akan ditangani oleh \textit{retry worker}.

\textit{Retry worker} dijalankan sebagai \textit{task} Tokio terpisah melalui
fungsi \texttt{spawn\_retry\_worker} yang dipanggil saat aplikasi diinisialisasi.
Worker berjalan dalam \textit{loop} tak terbatas yang setiap iterasinya membaca
seluruh entri dari antrian, mencoba mengirim ulang setiap entri yang belum
melampaui batas percobaan maksimum (\texttt{MAX\_ATTEMPTS\_BEFORE\_SKIP}), dan
menerapkan \textit{exponential backoff} pada jeda antar-iterasi. Jeda dihitung
sebagai $\min(\texttt{RETRY\_BASE\_SECS} \times 2^{\texttt{consecutive\_failures}},
\texttt{RETRY\_MAX\_SECS})$, sehingga semakin banyak kegagalan berturut-turut,
semakin panjang jeda sebelum percobaan berikutnya, tanpa tumbuh tak terbatas.
Entri yang telah melampaui batas percobaan tidak dikirim ulang tetapi tetap
disimpan dalam antrian untuk keperluan diagnostik.

\subsection{Implementasi Audit Receiver}

Komponen Audit Receiver diimplementasikan melalui \textit{handler}
\texttt{handle\_event} pada berkas \texttt{audit-trail/src/handlers.rs}
dan \textit{routing} pada \texttt{audit-trail/src/main.rs}. \textit{Endpoint}
\texttt{POST /api/events} menerima \textit{request} berisi objek
\texttt{EncryptedSignedEvent} dalam format JSON.

Proses pada \textit{handler} \texttt{handle\_event} terdiri atas tiga fase.
Fase pertama adalah verifikasi tanda tangan: ALS memverifikasi
\texttt{IotaSignature} yang ada di bidang \texttt{signature} terhadap
\textit{byte-byte ciphertext}, menggunakan \texttt{iota\_address} yang
disertakan sebagai referensi kunci publik. Apabila verifikasi gagal,
\textit{handler} mengembalikan respons \textit{error} tanpa memproses
\textit{event} lebih lanjut, sehingga \textit{ciphertext} yang dimanipulasi
pihak luar tidak pernah masuk ke tahap dekripsi.

Fase kedua adalah dekripsi \textit{payload}: ALS mendekripsi \texttt{enc\_aes\_key}
menggunakan kunci privat P-256 miliknya untuk memperoleh kunci sesi AES, kemudian
mendekripsi \texttt{ciphertext} menggunakan kunci sesi dan \texttt{nonce} tersebut
untuk memperoleh JSON \textit{plaintext} \texttt{AuditEvent}. AES-256-GCM
secara otomatis memverifikasi \textit{authentication tag} selama dekripsi; apabila
\textit{ciphertext} telah dirusak, dekripsi gagal dan \textit{event} ditolak.

Fase ketiga adalah pengiriman ke antrian internal: \textit{event} yang telah
terverifikasi dan terdekripsi diurai menjadi objek \texttt{AuditEvent} dan
diteruskan ke antrian asinkron \texttt{tokio::sync::mpsc} berkapasitas 10.000
\textit{event} melalui \texttt{state.audit\_tx.send(audit\_event).await}. Desain
ini memisahkan tanggung jawab penerimaan HTTP dari pemrosesan dan penulisan log,
sehingga ALS dapat menerima \textit{event} berikutnya tanpa menunggu penulisan
ke berkas log selesai.

Status implementasi komponen Audit Receiver ditunjukkan pada
Tabel~\ref{tab:status-implementasi-receiver}.

\begin{table}[htbp]
\centering
\caption{Status Implementasi Komponen Audit Receiver}
\label{tab:status-implementasi-receiver}
\renewcommand{\arraystretch}{1.3}
\begin{tabular}{|p{8cm}|p{3.5cm}|}
\hline
\textbf{Bagian} & \textbf{Status} \\
\hline
\textit{Endpoint} \texttt{POST /api/events} & Terimplementasi. \\
\hline
Verifikasi \texttt{IotaSignature} atas \textit{ciphertext} & Terimplementasi. \\
\hline
Dekripsi ECIES untuk memperoleh kunci sesi AES & Terimplementasi. \\
\hline
Dekripsi AES-256-GCM untuk memperoleh \textit{plaintext event} & Terimplementasi. \\
\hline
\textit{Parsing} \texttt{AuditEvent} dari \textit{plaintext} yang valid & Terimplementasi. \\
\hline
Pengiriman ke antrian asinkron internal (kapasitas 10.000) & Terimplementasi. \\
\hline
Penolakan \textit{event} dengan tanda tangan tidak valid & Terimplementasi. \\
\hline
\end{tabular}
\end{table}

\subsection{Implementasi Audit Logger}

Komponen Audit Logger diimplementasikan sebagai \texttt{struct AuditLogger}
pada berkas \texttt{audit-trail/src/audit.rs}. Komponen ini berjalan sebagai
\textit{task} Tokio terpisah yang di-\textit{spawn} pada saat aplikasi dimulai.

\texttt{AuditLogger} menyimpan dua bidang: \texttt{rx} bertipe
\texttt{Receiver<AuditEvent>} sebagai sisi penerima antrian asinkron, dan
\texttt{prev\_record\_hash} bertipe \texttt{Option<String>} sebagai \textit{hash}
dari \textit{record} sebelumnya untuk keperluan \textit{hash-chaining}.

Metode \texttt{run} menjalankan \textit{loop} utama menggunakan pola
\texttt{while let Some(event) = self.rx.recv().await}. Untuk setiap \textit{event},
\textit{loop} melakukan tiga langkah berurutan: (1) memanggil fungsi
\texttt{create\_audit\_record} untuk membentuk \texttt{AuditRecord} baru,
(2) memanggil \texttt{write\_audit\_record} untuk menulis \textit{record} ke
berkas log, dan (3) memperbarui \texttt{self.prev\_record\_hash} dengan
\textit{hash} dari \textit{record} yang baru saja ditulis.

Fungsi \texttt{create\_audit\_record} merealisasikan mekanisme
\textit{hash-chaining} dengan membangkitkan \texttt{record\_id} baru menggunakan
\texttt{Uuid::now\_v7()}, mengambil waktu saat ini menggunakan \texttt{Utc::now()},
kemudian mendelegasikan komputasi \textit{hash} ke
\texttt{Utils::calculate\_record\_hash}. Fungsi \texttt{calculate\_record\_hash}
menggabungkan \texttt{record\_id}, \texttt{timestamp}, \texttt{prev\_record\_hash},
dan isi \textit{event} menjadi satu objek JSON, kemudian menghitung SHA-256 dari
representasi JSON tersebut dan mengembalikannya sebagai \textit{string} heksadesimal.

Fungsi \texttt{write\_audit\_record} membuka berkas log dengan mode
\textit{append} menggunakan \texttt{OpenOptions::new().create(true).append(true)},
kemudian menulis representasi JSON dari \textit{record} diikuti karakter baris baru.
Penggunaan mode \textit{append} menjamin bahwa entri yang sudah ada tidak pernah
tertimpa, merealisasikan karakteristik \textit{append-only} yang disyaratkan oleh
KNF-TA-01. Konstanta \texttt{LOG\_FILE\_PATH} yang merujuk ke
\texttt{./logs/audit\_logs.log} didefinisikan pada \texttt{audit-trail/src/constants.rs}.

\subsection{Implementasi Log Rotation dan Archival Worker}

Komponen Log Rotation \& Archival Worker diimplementasikan sebagai fungsi
\texttt{Utils::spawn\_log\_rotation\_worker} pada berkas
\texttt{audit-trail/src/utils.rs}.

Worker berjalan sebagai \textit{task} Tokio terpisah yang berulang secara
periodik sesuai \texttt{LOG\_ROTATION\_INTERVAL\_SECS}. Pada setiap siklus,
worker melakukan langkah-langkah berikut secara berurutan.

\begin{enumerate}
    \item \textbf{Pembacaan dan validasi berkas log aktif.} Worker memeriksa
    apakah berkas \texttt{audit\_logs.log} ada dan tidak kosong. Apabila kosong,
    siklus dilewati.

    \item \textbf{Pembacaan seluruh \textit{record} untuk pembentukan
    \texttt{AuditBatch}.} Worker membaca seluruh baris berkas log untuk
    mengekstrak \texttt{first\_record\_id}, \texttt{last\_record\_id},
    \texttt{first\_record\_hash}, \texttt{final\_record\_hash},
    \texttt{record\_count}, \texttt{start\_time}, dan \texttt{end\_time}.

    \item \textbf{Penghitungan \textit{hash} berkas arsip.} Worker menghitung
    SHA-256 dari seluruh isi berkas log untuk menghasilkan \texttt{file\_hash}
    yang akan disimpan pada metadata \textit{on-chain}.

    \item \textbf{Penggantian berkas aktif.} Berkas log aktif diganti namanya
    menjadi berkas arsip bertanda waktu, dan berkas \texttt{audit\_logs.log}
    baru dibuat kosong untuk melanjutkan pencatatan.

    \item \textbf{Pengunggahan berkas arsip ke IPFS.} Berkas arsip diunggah
    ke IPFS Cluster menggunakan \textit{multipart upload} melalui \texttt{reqwest}
    dengan \texttt{BytesCodec}. CID yang dikembalikan oleh IPFS disimpan sebagai
    \texttt{ipfs\_cid} pada \texttt{AuditBatch}.

    \item \textbf{Publikasi metadata ke IOTA.} Worker memanggil fungsi
    \texttt{create\_log} pada \textit{smart contract} Move melalui Gas Station
    Server, menerbitkan \texttt{AuditBatch} yang telah terbentuk sebagai
    \textit{frozen object} pada jaringan IOTA.

    \item \textbf{Penghapusan berkas arsip lokal.} Setelah unggah IPFS dan
    publikasi IOTA berhasil, berkas arsip lokal dihapus untuk menghemat
    kapasitas penyimpanan.
\end{enumerate}

\subsection{Implementasi Smart Contract Move}

\textit{Smart contract} Move diimplementasikan pada modul
\texttt{als::audit\_log} dalam \textit{package} Move yang di-\textit{deploy} ke
jaringan IOTA Localnet. Modul ini menyediakan satu fungsi utama, yaitu
\texttt{create\_log}, yang menerima representasi JSON dari metadata rotasi dan
membuat objek \texttt{LogRecord} baru pada \textit{ledger} IOTA.

Aspek terpenting dari rancangan \textit{smart contract} ini adalah bahwa modul
\texttt{als::audit\_log} tidak menyediakan fungsi apa pun untuk mengubah atau
menghapus objek \texttt{LogRecord} yang telah dibuat. Satu-satunya operasi yang
tersedia adalah \texttt{create\_log}. Setiap objek \texttt{LogRecord} yang dibuat
langsung dibekukan melalui \texttt{transfer::freeze\_object}, sehingga tidak ada
entitas mana pun --- termasuk ALS itu sendiri --- yang dapat mengubah isinya
setelah transaksi pembuatannya \textit{final}. Rancangan ini secara langsung
mengatasi masalah M-TA-03 yang mengidentifikasi bahwa mekanisme pencatatan
eksisting pada DecMed masih menyediakan fungsi pembaruan pada \textit{smart
contract} Move yang ada.

\subsection{Implementasi Integrasi ALS ke DecMed}

Integrasi ALS ke dalam arsitektur DecMed dilakukan secara seragam pada seluruh
\textit{event source} yang telah diidentifikasi: PRE Server, \textit{Patient
Client}, \textit{Hospital Client}, dan \textit{Ministry Client}. Setiap
\textit{event source} diberikan penambahan modul \texttt{als} dengan pola
struktur yang konsisten di setiap komponen: berkas \texttt{als.rs} yang memuat
implementasi \texttt{ALSClient} dan \texttt{AlsQueue}, \texttt{types.rs} yang
mendefinisikan tipe-tipe data bersama, \texttt{crypto.rs} yang memuat fungsi
kriptografi, dan \texttt{queue.rs} yang mengimplementasikan \texttt{AlsQueue}
dan \textit{retry worker}.

\subsubsection{Tipe Data Bersama}

Pada setiap \textit{event source}, berkas \texttt{als/types.rs} mendefinisikan
ulang tipe-tipe data yang strukturnya setara dengan yang ada di ALS. Pendekatan
ini dipilih agar setiap komponen sumber \textit{event} tidak bergantung pada
\textit{crate} ALS secara langsung, mempertahankan independensi antarkomponen.
Varian \texttt{AuditEventDetails} yang diimplementasikan pada masing-masing
\textit{event source} dibatasi hanya pada jenis \textit{event} yang relevan
dengan aktivitas komponen tersebut; varian untuk \textit{event} yang berasal
dari komponen lain tidak didefinisikan.

\subsubsection{Implementasi ALSClient}

Setiap modul \texttt{als} menyediakan \textit{struct} \texttt{ALSClient} yang
mengelola seluruh alur enkapsulasi dan pengiriman \textit{event}. Berbeda dari
rancangan awal yang menggunakan kunci Ed25519 yang dibangkitkan secara acak,
implementasi akhir menggunakan \texttt{IotaKeyPair} yang sudah dimiliki oleh
aktor yang sedang aktif pada komponen. Kunci ini diambil dari \texttt{AppState}
komponen saat \textit{event} dibangkitkan, sehingga tanda tangan yang dihasilkan
secara langsung mencerminkan identitas aktor pada \textit{ledger} IOTA.

\texttt{ALSClient} menyediakan satu metode publik utama:
\texttt{send\_event\_from\_state}, yang menerima referensi \texttt{AppState},
objek \texttt{Event}, dan label \textit{event} bertipe \texttt{\&'static str}.
Metode ini mengambil \texttt{IotaKeyPair} dari \texttt{AppState}, membangun
\texttt{AuditEvent} dengan \texttt{source\_component} yang sesuai komponen dan
\texttt{source\_timestamp} berupa \texttt{Utc::now()}, kemudian men-\textit{spawn}
\textit{task} Tokio yang memanggil \texttt{build\_and\_send} secara asinkron.
Pola \textit{spawn-and-forget} ini memastikan bahwa operasi audit tidak
menambah latensi pada alur layanan utama.

\textit{Retry worker} dijalankan sekali saat \textit{startup} melalui
\texttt{ALSClient::start\_retry\_worker()}, yang pada aplikasi Tauri
menggunakan \texttt{tauri::async\_runtime::spawn} dan pada PRE Server
menggunakan \texttt{tokio::spawn}.

\subsubsection{Integrasi ke State Komponen}

Pada PRE Server, \texttt{ALSClient} tidak disimpan sebagai bidang tersendiri
pada \texttt{AppState} karena semua dependensi yang dibutuhkan (\texttt{IotaKeyPair}
dan \texttt{iota\_address}) sudah tersedia pada \texttt{AppState}. Pada aplikasi
\textit{client} berbasis Tauri, \texttt{AppState} berisi \texttt{keys\_entry}
terenkripsi yang dibuka menggunakan PIN sesi aktif pada saat \textit{event}
dibangkitkan.

\subsubsection{Distribusi Event per Event Source}

\textit{Patient Client}, \textit{Hospital Client}, dan \textit{Ministry Client}
merupakan aplikasi \textit{client} berbasis Tauri. Pada ketiga aplikasi tersebut,
modul \texttt{als} diimplementasikan pada sisi \textit{backend} Rust dari aplikasi
Tauri, sehingga enkripsi, penandatanganan, dan pengiriman \textit{event} tidak
terekspos ke lapisan antarmuka SvelteKit/TypeScript. Distribusi \textit{event}
berdasarkan \textit{event source} ditunjukkan pada Tabel~\ref{tab:event-source-distribution}.

Perlu dicatat bahwa \textit{event} EV8 (IOTA Transaction Submission) dan EV9 (Gas Sponsorship Request) tercatat pada sisi pemohon, bukan pada Gas Station
Server.